\section*{الفصل \ref{ch:g8:colors-spectra} --- ألوان الضوء وأطيافه}

\begin{solution}{exo:g8:colors-spectra:1}
الشريط المتّصل من الألوان المختبئة في \omterm{def:g1:light-and-shadows:source}{الضوء} الأبيض، مفكوكًا؛
والتبديد هو الفكّ --- إذ ينكسر كل لون بمقداره الخاصّ. والأحمر
أقلّ انكسارًا، والبنفسجي أكثر.
\end{solution}

\begin{solution}{exo:g8:colors-spectra:2}
قوس قزح يتكوّن في نصف السماء المقابل للشمس --- \omterm{def:g1:light-and-shadows:source}{فضوء} الشمس يدخل
قطرات المطر من خلف الصائد. وكل قطرة مطر تؤدّي دور الموشور (مع
ارتداد \omterm{def:g5:mirrors-reflection:reflection}{مرآة} في داخلها).
\end{solution}

\begin{solution}{exo:g8:colors-spectra:3}
الفكّ: الموشور يفتح \omterm{def:g1:light-and-shadows:source}{الضوء} الأبيض مروحةً في الطيف. وإعادة
التعبئة: القرص الدوّار ذو ألوان قوس قزح يتشوّش إلى بياض ---
والموشور الثاني يطوي \omterm{prop:g8:colors-spectra:mixture}{طيفًا} ثانية في \omterm{def:g6:light-propagation:ray}{حزمة} بيضاء واحدة.
\end{solution}

\begin{solution}{exo:g8:colors-spectra:4}
العشب يمتصّ معظم المزيج ويعيد الأخضر أساسًا. والصفحة تعيد كل
شيء تقريبًا --- فبيضاء؛ والسبّورة لا تعيد شيئًا تقريبًا بل
تمتصّ الجميع --- فسوداء.
\end{solution}

\begin{solution}{exo:g8:colors-spectra:5}
\omterm{def:g1:light-and-shadows:source}{الضوء} الأبيض: الوشاح يعيد أصفره --- وشاح أصفر. \omterm{def:g1:light-and-shadows:source}{والضوء} الأزرق:
يلقى العرض (الأزرق) الجوابَ (الأصفر) في تقاطع خالٍ --- فيبدو
الوشاح أسود.
\end{solution}

\begin{solution}{exo:g8:colors-spectra:6}
التفاحة لا تستطيع أن تعيد إلا الأحمر، والكشّاف الأزرق لا يجيء
بشيء منه: فسوداء. والمنديل يعيد ما يصل --- وقد عُرض عليه
الأزرق فيعيد الأزرق: منديل أزرق.
\end{solution}

\begin{solution}{exo:g8:colors-spectra:7}
الأحمر والأخضر والأزرق. الأحمر $+$ الأخضر: أصفر. والأخضر $+$
الأزرق: سماوي. والأزرق $+$ الأحمر: أرجواني. والثلاثة جميعًا:
أبيض.
\end{solution}

\begin{solution}{exo:g8:colors-spectra:8}
شرائط دقيقة حمراء وخضراء وزرقاء لا غير، كلها مضاءة --- ولا
بياض في أيّ موضع. والمزج يقع في عينك، وهي أبعد من أن تفصل
الشرائط.
\end{solution}

\begin{solution}{exo:g8:colors-spectra:9}
الأضواء تجمع ما تجيء به، ومجموعة الألوان كاملةً معًا بيضاء ---
فالجمع يتمّ المزيج. والأصباغ تمتصّ، فيطرح كلٌّ حصته من \omterm{def:g1:light-and-shadows:source}{الضوء}،
ومجموعة الطروح كاملةً لا تُبقي شيئًا يُذكر --- ظلمة بقرار
لجنة.
\end{solution}

\begin{solution}{exo:g8:colors-spectra:10}
\omterm{def:g2:air-around-us:air}{الهواء} يسرق الأزرق تفضيلًا؛ والسماء كلها تتلقّى الأزرق المسروق
وترسله إلينا قبّةً زرقاء نهارية؛ والباقي الأحمر البرتقالي ينجو
من طريق \omterm{def:g1:day-and-night:daynight}{الغروب} الملامس. أما \omterm{def:g2:moon-in-the-sky:moon}{القمر} المخسوف فلا يضيئه إلا \omterm{def:g1:light-and-shadows:source}{ضوء}
لامس حافة الأرض --- مرشَّحًا \omterm{def:g1:day-and-night:daynight}{بالغروب} حول \omterm{def:g4:solar-system:planet}{الكوكب} كله --- فيتوهّج
بالناجي الأحمر نفسه.
\end{solution}

\begin{solution}{exo:g8:colors-spectra:11}
المصابيح الدافئة المائلة إلى الحمرة لم تعرض إلا القليل من
الأخضر؛ وجواب السترة الحقيقي --- أخضر في معظمه مع بعض الأحمر
يصنعان مزيجًا كدرًا --- لم يستطع أن يُظهر إلا حصته الصديقة
للأحمر، فقُرئت بنّية دافئة. أما عرض \omterm{def:g1:light-and-shadows:source}{ضوء} النهار الكامل فأتاح
للجواب الأخضر أن يتكلّم: خديعة \omterm{def:g1:light-and-shadows:source}{بضوء} بخيل، تدينها قاعدة
التقاطع.
\end{solution}

\begin{solution}{exo:g8:colors-spectra:12}
الحروف زرقاء، والأرضية سوداء. فتحت أضواء الصالة البيضاء تعيد
الحروف الزرقاء زرقتها بينما لا تعيد الأرضية السوداء شيئًا:
حروف زرقاء ساطعة على سواد --- مقروءة. وتحت الغمر الأحمر العميق
لا يُعرَض على الحروف الزرقاء شيء تستطيع إعادته، ولا تعيد
الأرضية السوداء شيئًا كعادتها: سواد منتظم، واللافتة مختفية.
أما الحروف البيضاء فكانت ستخون الختام: فالمُعيد الشامل يجيب
حتى الغمر الأحمر بحروف حمراء ساطعة، مقروءة على الأرضية
السوداء.
\end{solution}

\begin{solution}{pb:g8:colors-spectra:1}
\textbf{1.} أبيض --- فالألوان الأساسية الثلاثة بكامل قوّتها
تجتمع في المزيج الكامل.
\textbf{2.} أصفر.
\textbf{3.} القميص الأبيض: يعيد العرض --- فأصفر. والوشاح
الأزرق: لم يُعرض عليه أزرق --- فأسود. والثوب الأصفر: جوابه
الأصفر يطابق العرض الأصفر --- فأصفر باهر.
\textbf{4.} خلفية زرقاء. فالقماش الأبيض يعيد ما يرسله
المصمّم: خلفية واحدة، وكل لون عند الطلب --- نظير الشاشة
البيضاء في السينما.
\textbf{5.} معطف أخضر: فتحت الغمر الأحمر لا يلقى جوابه الأخضر
أخضر --- فشرّير أسود؛ وتحت \omterm{def:g1:light-and-shadows:source}{الضوء} الأبيض يتيح العرض الكامل
للأخضر أن يتكلّم --- فأخضر غنيّ.
\textbf{6.} رداء أبيض --- المُعيد الشامل: أحمر تحت الكشّاف
الأحمر، وأخضر تحت الأخضر، وأزرق تحت الأزرق، وأبيض تحت الثلاثة
معًا؛ وهو دائمًا أسطع شيء على الخشبة، ودائمًا ``يبدو أبيض''
\omterm{def:g6:measurement-in-science:measuring}{بالقياس} إلى المشهد.
\textbf{7.} أيّ إضاءة \emph{بلا} أحمر: الكشّاف الأزرق، أو
الأخضر، أو كلاهما معًا --- فإذ لا يُعرض عليه أحمر، لا يعيد
الوشاح شيئًا ويذوب في الستارة السوداء. ولحظة ينضمّ الكشّاف
الأحمر يصير العرض حاويًا بالضبط ما يجيب به الوشاح: فيشتعل
أحمر، وتموت الحيلة.
\textbf{8.} الذهبي بوصفه أصفر دافئًا يجيب الأحمر والأخضر
بسخاء، والأزرق بالكاد. وضبط المساء الغنيّ بالأزرق لم يعرض
عليه إلا القليل مما يستطيع إعادته: فحروف كدرة خافتة. أما مزيج
الظهيرة الكامل فأتاح للأصفر أن يشتعل.
\textbf{9.} الخطّ الأحمر: أحمر باهر. والبرتقالي والأصفر:
نسختان خافتتان مائلتان إلى الحمرة من نفسيهما (فجوابهما يحوي
بعض الأحمر). والأخضر: أسود. والأزرق والبنفسجي: أسودان
(والبنفسجي بالكاد جمرة داكنة). فينهار القوس إلى خراب أحمر
وأسود.
\textbf{10.} الخطّان الأحمر والأزرق يحييان؛ والبنفسجي يتوهّج
خافتًا (فهو يجيب الأزرق وبعض الأحمر)؛ والبرتقالي والأصفر لا
يُظهران إلا حصّتيهما الحمراوين الضعيفتين؛ ويبقى الخطّ الأخضر
أسود --- \omterm{def:g1:light-and-shadows:source}{فالضوء} الأرجواني لا يحوي أخضر.
\textbf{11.} \omterm{def:g1:light-and-shadows:source}{الضوء} الأبيض وحده --- المزيج الكامل --- يُظهر
الخطوط الستّة كلها: فكل خطّ يعيد لونه هو، ولا يجيب عن الستّة
جميعًا إلا عرض يحوي الستّة جميعًا. أما الكشّاف الأخضر، مهما
سطع، فيضيء خطًّا واحدًا ويقتل خمسة.
\textbf{12.} مثلًا: ``الأضواء الممزوجة \emph{تجمع}: فالأحمر
والأخضر والأزرق تبني كل لون، وتبني معًا الأبيض. وكل جسم
\emph{يطرح}: فلا يعيد إلا حصته من الذي جاء به \omterm{def:g1:light-and-shadows:source}{الضوء} --- فاعرض
عليه ما لا يستطيع إعادته يلبس السواد.''
\end{solution}
